\chapter{Code}

Full Framework Codebase can be found here: https://github.com/vbalab/FITS

% \begin{lstlisting}[language=Python, caption={Real-time AO region detection code functionality}]
% ================================================
% FILE: src/fits/modelling/framework.py
% ================================================
% import pickle
% from typing import Optional
% from datetime import datetime
% from dataclasses import dataclass
% from abc import ABC, abstractmethod

% import torch
% import numpy as np
% import matplotlib.pyplot as plt
% from torch import nn
% from tqdm import tqdm
% from torch.utils.data import DataLoader
% from IPython.display import clear_output

% from fits.config import TRAINING_PATH, EVALUATION_PATH
% from fits.dataframes.dataset import ForecastingData, NormalizationStats


% @dataclass
% class ForecastedData:
%     # `nsample` - number of Monte-Carlo forecast samples
%     forecasted_data: torch.Tensor  # [B, nsample, L, K] float
%     forecast_mask: torch.Tensor  # [B, L, K] int
%     # 0 - context (history)
%     # 1 - horizon to forecast (to be generated)
%     observed_data: torch.Tensor  # [B, L, K] float
%     observed_mask: torch.Tensor  # [B, L, K] int
%     # 0 - not observed
%     # 1 - observed (ground truth)
%     time_points: torch.Tensor  # [B, L] float


% @dataclass
% class ModelConfig:
%     """Base model configuration used by :class:`ForecastingModel`.

%     The configuration is intentionally typed to avoid passing around loose
%     dictionaries.
%     """

%     name: Optional[str] = None


% class ForecastingModel(nn.Module, ABC):
%     """Base class for all forecasting models in the project."""

%     def __init__(self, config: ModelConfig):
%         super().__init__()
%         self.config = config

%         if not torch.cuda.is_available():
%             raise ValueError("CUDA is not available")

%         self.device = torch.device("cuda")

%     @property
%     def model_name(self) -> str:
%         return self.config.name or self.__class__.__name__

%     @abstractmethod
%     def forward(self, batch: ForecastingData):
%         """Compute the training loss for a batch."""

%     @abstractmethod
%     def evaluate(self, batch: ForecastingData, n_samples: int) -> ForecastedData:
%         """Run model-specific evaluation and return generated samples."""


% @dataclass
% class EMA:
%     """Exponential Moving Average of model parameters."""

%     decay: float = 0.999
%     device: Optional[torch.device] = (
%         None  # e.g. torch.device("cpu") to store EMA on CPU
%     )

%     def __post_init__(self):
%         self.shadow = {}
%         self._backup = {}

%     @torch.no_grad()
%     def register(self, model: torch.nn.Module) -> None:
%         self.shadow = {}
%         for name, p in model.named_parameters():
%             if p.requires_grad:
%                 self.shadow[name] = (
%                     p.detach().clone().to(self.device)
%                     if self.device
%                     else p.detach().clone()
%                 )

%     @torch.no_grad()
%     def update(self, model: torch.nn.Module) -> None:
%         d = self.decay
%         for name, p in model.named_parameters():
%             if not p.requires_grad:
%                 continue
%             assert (
%                 name in self.shadow
%             ), "EMA.register(model) must be called before EMA.update(model)."
%             new = p.detach()
%             if self.device:
%                 new = new.to(self.device)
%             self.shadow[name].mul_(d).add_(new, alpha=(1.0 - d))

%     @torch.no_grad()
%     def apply(self, model: torch.nn.Module) -> None:
%         """Swap model params to EMA params (stores current params for later restore)."""
%         self._backup = {}
%         for name, p in model.named_parameters():
%             if not p.requires_grad:
%                 continue
%             self._backup[name] = p.detach().clone()
%             ema_p = self.shadow[name]
%             if ema_p.device != p.device:
%                 ema_p = ema_p.to(p.device)
%             p.copy_(ema_p)

%     @torch.no_grad()
%     def restore(self, model: torch.nn.Module) -> None:
%         """Restore params saved by apply()."""
%         for name, p in model.named_parameters():
%             if not p.requires_grad:
%                 continue
%             p.copy_(self._backup[name])
%         self._backup = {}


% def Train(
%     model: ForecastingModel,
%     train_loader: DataLoader,
%     valid_loader: DataLoader,
%     lr: float = 1.0e-3,
%     epochs: int = 500,
%     valid_epoch_interval: int = 20,
%     verbose: bool = True,
%     warmup_epochs: int = 10,
%     warmup_start_factor: float = 0.1,  # initial lr = lr * warmup_start_factor
%     grad_clip_norm: float | None = 0.3,
%     weight_decay: float = 1e-6,
%     use_ema: bool = False,
%     ema_decay: float = 0.995,  # if use_ema=True
%     ema_eval: bool = True,  # if use_ema=True
%     ema_save: bool = True,  # if use_ema=True
%     folder_name: str | None = None,
% ):
%     if not folder_name:
%         current_time = datetime.now().strftime("%Y%m%d_%H%M%S")
%         folder_name = f"{model.model_name}_{current_time}"

%     folder_path = TRAINING_PATH / folder_name
%     folder_path.mkdir(parents=True, exist_ok=True)

%     opt = torch.optim.Adam(model.parameters(), lr=lr, weight_decay=weight_decay)

%     # first 75% epoechs: 1e-3
%     # then  15% epoechs: 1e-3 * 0.1 = 1e-4
%     # then  10% epoechs: 1e-4 * 0.1 = 1e-5
%     p1 = int(0.75 * epochs)
%     p2 = int(0.9 * epochs)

%     scheduler: (
%         torch.optim.lr_scheduler.LambdaLR | torch.optim.lr_scheduler.MultiStepLR | None
%     ) = None
%     if warmup_epochs > 0:

%         def lr_lambda(epoch: int) -> float:
%             warm = min(1.0, (epoch + 1) / float(warmup_epochs))
%             # approximate MultiStepLR behavior with epoch-based factor:
%             # after passing milestones, apply gamma^k
%             k = 0
%             if (epoch + 1) > p1:
%                 k += 1
%             if (epoch + 1) > p2:
%                 k += 1
%             return (warmup_start_factor + (1.0 - warmup_start_factor) * warm) * (0.1**k)

%         scheduler = torch.optim.lr_scheduler.LambdaLR(opt, lr_lambda=lr_lambda)
%     else:
%         scheduler = torch.optim.lr_scheduler.MultiStepLR(
%             opt, milestones=[p1, p2], gamma=0.1
%         )

%     ema: EMA | None = None
%     if use_ema:
%         ema = EMA(decay=ema_decay)
%         ema.register(model)

%     metrics: dict[str, list[tuple[int, float]]] = {"train_loss": [], "test_loss": []}
%     best_valid_loss = float("inf")

%     for epoch_no in range(epochs):
%         epoch_loss = 0.0
%         model.train()

%         with tqdm(train_loader) as it:
%             for batch_no, train_batch in enumerate(it, start=1):
%                 opt.zero_grad()

%                 loss = model(train_batch)
%                 epoch_loss += loss.item()

%                 loss.backward()
%                 if grad_clip_norm is not None:
%                     nn.utils.clip_grad_norm_(model.parameters(), grad_clip_norm)
%                 opt.step()

%                 if ema:
%                     ema.update(model)

%                 it.set_postfix(
%                     ordered_dict={
%                         "avg_epoch_loss": epoch_loss / batch_no,
%                         "epoch": epoch_no,
%                     },
%                     refresh=False,
%                 )

%         scheduler.step()
%         metrics["train_loss"].append((epoch_no + 1, epoch_loss))

%         if (epoch_no + 1) % valid_epoch_interval == 0:
%             model.eval()

%             if ema and ema_eval:
%                 ema.apply(model)

%             valid_loss = 0.0
%             valid_batches = 0

%             with torch.no_grad():
%                 with tqdm(valid_loader) as it:
%                     for batch_no, valid_batch in enumerate(it, start=1):
%                         loss = model(valid_batch)
%                         valid_loss += loss.item()
%                         valid_batches = batch_no

%                         it.set_postfix(
%                             ordered_dict={
%                                 "valid_avg_epoch_loss": valid_loss / batch_no,
%                                 "epoch": epoch_no,
%                             },
%                             refresh=False,
%                         )

%             metrics["test_loss"].append((epoch_no + 1, valid_loss))

%             avg_valid_loss = valid_loss / max(valid_batches, 1)
%             if best_valid_loss > avg_valid_loss:
%                 torch.save(model.state_dict(), folder_path / "best_model.pth")

%                 best_valid_loss = avg_valid_loss
%                 print("\n best loss is updated to ", avg_valid_loss, "at", epoch_no)

%             if ema and ema_eval:
%                 ema.restore(model)

%         if verbose:
%             clear_output(True)
%             plt.figure(figsize=(12, 4))

%             for i, (name, history) in enumerate(sorted(metrics.items())):
%                 plt.subplot(1, len(metrics), i + 1)
%                 plt.title(name)
%                 plt.plot(*zip(*history))
%                 plt.grid()

%             plt.gca().text(
%                 0.02,
%                 0.98,
%                 f"lr = {opt.param_groups[0]['lr']:.2e}",
%                 transform=plt.gca().transAxes,
%                 va="top",
%                 ha="left",
%             )

%             plt.savefig(folder_path / "training.png")
%             plt.show()

%         if ema and ema_save:
%             ema.apply(model)

%         torch.save(model.state_dict(), folder_path / "model.pth")

%         if ema and ema_save:
%             ema.restore(model)


% def CalcQuantileLoss(target, forecast, q: float, eval_points) -> torch.Tensor:
%     return 2 * torch.sum(
%         torch.abs((forecast - target) * eval_points * ((target <= forecast) * 1.0 - q))
%     )


% def CalcQuantileCRPS(target, forecast, eval_points, center, scale):
%     target = target * scale + center
%     forecast = forecast * scale + center

%     quantiles = np.arange(0.05, 1.0, 0.05)
%     denom = torch.sum(torch.abs(target * eval_points))
%     CRPS = 0
%     for i in range(len(quantiles)):
%         q_pred = []
%         for j in range(len(forecast)):
%             q_pred.append(torch.quantile(forecast[j : j + 1], quantiles[i], dim=1))
%         q_pred = torch.cat(q_pred, 0)
%         q_loss = CalcQuantileLoss(target, q_pred, quantiles[i], eval_points)
%         CRPS += q_loss / denom
%     return CRPS.item() / len(quantiles)


% def CalcQuantileCRPSSum(target, forecast, eval_points, center, scale):
%     eval_points = eval_points.float().mean(-1)
%     target = target * scale + center
%     target = target.sum(-1)
%     forecast = forecast * scale + center

%     quantiles = np.arange(0.05, 1.0, 0.05)
%     denom = torch.sum(torch.abs(target * eval_points))
%     CRPS = 0
%     for i in range(len(quantiles)):
%         q_pred = torch.quantile(forecast.sum(-1), quantiles[i], dim=1)
%         q_loss = CalcQuantileLoss(target, q_pred, quantiles[i], eval_points)
%         CRPS += q_loss / denom
%     return CRPS.item() / len(quantiles)


% @torch.no_grad()
% def Evaluate(
%     model: ForecastingModel,
%     test_loader: DataLoader,
%     normalization: NormalizationStats,
%     nsample: int = 5,
%     folder_name: str | None = None,
% ):
%     if not folder_name:
%         current_time = datetime.now().strftime("%Y%m%d_%H%M%S")
%         folder_name = f"{model.model_name}_{current_time}"

%     folder_path = EVALUATION_PATH / folder_name
%     folder_path.mkdir(parents=True, exist_ok=True)

%     model.eval()
%     mse_total = 0.0
%     mae_total = 0.0
%     evalpoints_total = 0.0

%     min_tensor = torch.as_tensor(normalization.min, device=model.device)
%     max_tensor = torch.as_tensor(normalization.max, device=model.device)
%     center_tensor = (max_tensor + min_tensor) / 2
%     scale_tensor = (max_tensor - min_tensor) / 2
%     scale_tensor = torch.where(
%         scale_tensor == 0, torch.ones_like(scale_tensor), scale_tensor
%     )

%     all_forecasted_data: list[torch.Tensor] = []
%     all_forecast_mask: list[torch.Tensor] = []
%     all_observed_data: list[torch.Tensor] = []
%     all_observed_mask: list[torch.Tensor] = []
%     all_time_points: list[torch.Tensor] = []

%     with tqdm(test_loader, mininterval=5.0, maxinterval=50.0) as it:
%         for batch_no, test_batch in enumerate(it, start=1):
%             f: ForecastedData = model.evaluate(test_batch, nsample)

%             all_forecasted_data.append(f.forecasted_data)
%             all_forecast_mask.append(f.forecast_mask)
%             all_observed_data.append(f.observed_data)
%             all_observed_mask.append(f.observed_mask)
%             all_time_points.append(f.time_points)

%             forecasted_data_median = f.forecasted_data.median(dim=1).values

%             mse_current = (
%                 ((forecasted_data_median - f.observed_data) * f.forecast_mask) ** 2
%             ) * (scale_tensor**2)
%             mae_current = (
%                 torch.abs((forecasted_data_median - f.observed_data) * f.forecast_mask)
%             ) * scale_tensor

%             mse_total += mse_current.sum().item()
%             mae_total += mae_current.sum().item()
%             evalpoints_total += f.forecast_mask.sum().item()

%             it.set_postfix(
%                 ordered_dict={
%                     "rmse_total": np.sqrt(mse_total / evalpoints_total),
%                     "mae_total": mae_total / evalpoints_total,
%                     "batch_no": batch_no,
%                 },
%                 refresh=True,
%             )

%     all_forecasted_data_tensor = torch.cat(all_forecasted_data, dim=0)
%     all_forecast_mask_tensor = torch.cat(all_forecast_mask, dim=0)
%     all_observed_data_tensor = torch.cat(all_observed_data, dim=0)
%     all_observed_mask_tensor = torch.cat(all_observed_mask, dim=0)
%     all_time_points_tensor = torch.cat(all_time_points, dim=0)

%     with open(folder_path / f"generated_outputs_nsample{nsample}.pk", "wb") as file:
%         pickle.dump(
%             [
%                 all_forecasted_data_tensor,
%                 all_forecast_mask_tensor,
%                 all_observed_data_tensor,
%                 all_observed_mask_tensor,
%                 all_time_points_tensor,
%                 scale_tensor.cpu(),
%                 center_tensor.cpu(),
%             ],
%             file,
%         )

%     crps = CalcQuantileCRPS(
%         all_observed_data_tensor,
%         all_forecasted_data_tensor,
%         all_forecast_mask_tensor,
%         center_tensor,
%         scale_tensor,
%     )
%     crps_sum = CalcQuantileCRPSSum(
%         all_observed_data_tensor,
%         all_forecasted_data_tensor,
%         all_forecast_mask_tensor,
%         center_tensor,
%         scale_tensor,
%     )

%     rmse = np.sqrt(mse_total / evalpoints_total)
%     mae = mae_total / evalpoints_total

%     metrics = [rmse, mae, crps, crps_sum]
%     metrics = [float(metric) for metric in metrics]
%     with open(folder_path / f"result_nsample{nsample}.pk", "wb") as file:
%         pickle.dump(metrics, file)

%     print("RMSE:", rmse)
%     print("MAE:", mae)
%     print("CRPS:", crps)
%     print("CRPS_sum:", crps_sum)



% ================================================
% FILE: src/fits/modelling/FITSJ/decomposition.py
% ================================================
% import pickle
% from datetime import datetime
% from pathlib import Path
% from typing import Sequence

% import torch
% import matplotlib.pyplot as plt
% from torch.utils.data import DataLoader
% from tqdm import tqdm

% from fits.config import EVALUATION_PATH
% from fits.dataframes.dataset import NormalizationStats
% from fits.modelling.FITSJ.transformer import Decomposition
% from fits.modelling.FITSJ.model import FITSModel


% def _normalize_feature_indices(
%     feature_index: int | Sequence[int] | None,
%     n_features: int,
% ) -> list[int]:
%     if feature_index is None:
%         return list(range(n_features))
%     if isinstance(feature_index, int):
%         feature_index = [feature_index]
%     return [idx for idx in feature_index if 0 <= idx < n_features]


% @torch.no_grad()
% def EvaluateFITSJWithDecomposition(
%     model: FITSModel,
%     test_loader: DataLoader,
%     normalization: NormalizationStats,
%     nsample: int = 5,
%     folder_name: str | None = None,
% ):
%     if not folder_name:
%         current_time = datetime.now().strftime("%Y%m%d_%H%M%S")
%         folder_name = f"{model.model_name}_{current_time}"

%     folder_path = EVALUATION_PATH / folder_name
%     folder_path.mkdir(parents=True, exist_ok=True)

%     model.eval()

%     min_tensor = torch.as_tensor(normalization.min, device=model.device)
%     max_tensor = torch.as_tensor(normalization.max, device=model.device)
%     center_tensor = (max_tensor + min_tensor) / 2
%     scale_tensor = (max_tensor - min_tensor) / 2
%     scale_tensor = torch.where(
%         scale_tensor == 0, torch.ones_like(scale_tensor), scale_tensor
%     )

%     all_forecasted_data: list[torch.Tensor] = []
%     all_forecast_mask: list[torch.Tensor] = []
%     all_observed_data: list[torch.Tensor] = []
%     all_observed_mask: list[torch.Tensor] = []
%     all_time_points: list[torch.Tensor] = []
%     all_trend: list[torch.Tensor] = []
%     all_season: list[torch.Tensor] = []
%     all_jump: list[torch.Tensor] = []
%     all_error: list[torch.Tensor] = []

%     with tqdm(test_loader, mininterval=5.0, maxinterval=50.0) as it:
%         for test_batch in it:
%             forecasted, decomposed = model.decomposition_evaluate(test_batch, nsample)

%             all_forecasted_data.append(forecasted.forecasted_data)
%             all_forecast_mask.append(forecasted.forecast_mask)
%             all_observed_data.append(forecasted.observed_data)
%             all_observed_mask.append(forecasted.observed_mask)
%             all_time_points.append(forecasted.time_points)

%             all_trend.append(decomposed.trend)
%             all_season.append(decomposed.season)
%             all_jump.append(decomposed.jump_term)
%             all_error.append(decomposed.error)

%     all_forecasted_data_tensor = torch.cat(all_forecasted_data, dim=0)
%     all_forecast_mask_tensor = torch.cat(all_forecast_mask, dim=0)
%     all_observed_data_tensor = torch.cat(all_observed_data, dim=0)
%     all_observed_mask_tensor = torch.cat(all_observed_mask, dim=0)
%     all_time_points_tensor = torch.cat(all_time_points, dim=0)

%     with open(folder_path / f"generated_outputs_nsample{nsample}.pk", "wb") as file:
%         pickle.dump(
%             [
%                 all_forecasted_data_tensor,
%                 all_forecast_mask_tensor,
%                 all_observed_data_tensor,
%                 all_observed_mask_tensor,
%                 all_time_points_tensor,
%                 scale_tensor.cpu(),
%                 center_tensor.cpu(),
%             ],
%             file,
%         )

%     decomposed = Decomposition(
%         trend=torch.cat(all_trend, dim=0),
%         season=torch.cat(all_season, dim=0),
%         jump_term=torch.cat(all_jump, dim=0),
%         error=torch.cat(all_error, dim=0),
%     )
%     with open(folder_path / f"decomposition_outputs_nsample{nsample}.pk", "wb") as file:
%         pickle.dump(decomposed, file)


% def PlotFITSJDecomposition(
%     eval_foldername: str | Path,
%     nsample: int = 5,
%     sample_index: int = 0,
%     feature_index: int | Sequence[int] | None = None,
%     figsize: tuple[int, int] = (12, 10),
% ) -> None:
%     evaluation_dir = Path(f"../data/models/evaluation/{eval_foldername}")
%     generated_path = evaluation_dir / f"generated_outputs_nsample{nsample}.pk"
%     decomposition_path = evaluation_dir / f"decomposition_outputs_nsample{nsample}.pk"

%     with open(generated_path, "rb") as f:
%         (
%             forecasted_data,
%             forecast_mask,
%             observed_data,
%             observed_mask,
%             time_points,
%             scale_tensor,
%             center_tensor,
%         ) = pickle.load(f)

%     with open(decomposition_path, "rb") as f:
%         decomposed: Decomposition = pickle.load(f)

%     forecasted_data = forecasted_data.cpu()
%     forecast_mask = forecast_mask.cpu()
%     observed_data = observed_data.cpu()
%     observed_mask = observed_mask.cpu()
%     time_points = time_points.cpu()
%     scale_tensor = scale_tensor.cpu()
%     center_tensor = center_tensor.cpu()

%     n_features = observed_data.shape[-1]
%     selected_features = _normalize_feature_indices(feature_index, n_features)
%     if not selected_features:
%         raise ValueError("No valid feature indices provided.")

%     ncols = len(selected_features)
%     nrows = 5
%     fig, axes = plt.subplots(nrows=nrows, ncols=ncols, figsize=figsize, sharex=True)
%     if ncols == 1:
%         axes = axes.reshape(nrows, 1)

%     time_axis = time_points[sample_index].numpy()

%     forecast_samples = forecasted_data[sample_index]
%     forecast_median = forecast_samples.median(dim=0).values
%     forecast_lower, forecast_upper = torch.quantile(
%         forecast_samples, torch.tensor([0.1, 0.9]), dim=0
%     )

%     trend_samples = decomposed.trend[sample_index].cpu()
%     season_samples = decomposed.season[sample_index].cpu()
%     jump_samples = decomposed.jump_term[sample_index].cpu()
%     error_samples = decomposed.error[sample_index].cpu()

%     trend_median = trend_samples.median(dim=0).values
%     season_median = season_samples.median(dim=0).values
%     jump_median = jump_samples.median(dim=0).values
%     error_median = error_samples.median(dim=0).values

%     trend_lower, trend_upper = torch.quantile(
%         trend_samples, torch.tensor([0.1, 0.9]), dim=0
%     )
%     season_lower, season_upper = torch.quantile(
%         season_samples, torch.tensor([0.1, 0.9]), dim=0
%     )
%     jump_lower, jump_upper = torch.quantile(
%         jump_samples, torch.tensor([0.1, 0.9]), dim=0
%     )
%     error_lower, error_upper = torch.quantile(
%         error_samples, torch.tensor([0.1, 0.9]), dim=0
%     )

%     for col_idx, feat in enumerate(selected_features):
%         horizon_mask = forecast_mask[sample_index, :, feat].bool()
%         observed_series = observed_data[sample_index, :, feat]
%         observed_series_mask = observed_mask[sample_index, :, feat].bool()
%         observed_horizon_mask = horizon_mask & observed_series_mask

%         scale = scale_tensor[feat]
%         center = center_tensor[feat]

%         observed_denorm = observed_series * scale + center
%         forecast_denorm = forecast_median[:, feat] * scale + center
%         forecast_lower_denorm = forecast_lower[:, feat] * scale + center
%         forecast_upper_denorm = forecast_upper[:, feat] * scale + center

%         trend_denorm = trend_median[:, feat] * scale + center
%         trend_lower_denorm = trend_lower[:, feat] * scale + center
%         trend_upper_denorm = trend_upper[:, feat] * scale + center

%         season_denorm = season_median[:, feat] * scale
%         season_lower_denorm = season_lower[:, feat] * scale
%         season_upper_denorm = season_upper[:, feat] * scale

%         jump_denorm = jump_median[:, feat] * scale
%         jump_lower_denorm = jump_lower[:, feat] * scale
%         jump_upper_denorm = jump_upper[:, feat] * scale

%         error_denorm = error_median[:, feat] * scale
%         error_lower_denorm = error_lower[:, feat] * scale
%         error_upper_denorm = error_upper[:, feat] * scale

%         axes[0, col_idx].plot(
%             time_axis[observed_horizon_mask.numpy()],
%             observed_denorm[observed_horizon_mask].numpy(),
%             color="black",
%             label="Observed",
%         )
%         axes[0, col_idx].plot(
%             time_axis[horizon_mask.numpy()],
%             forecast_denorm[horizon_mask].numpy(),
%             color="tab:green",
%             label="Forecast (median)",
%         )
%         axes[0, col_idx].fill_between(
%             time_axis,
%             forecast_lower_denorm.numpy(),
%             forecast_upper_denorm.numpy(),
%             where=horizon_mask.numpy(),
%             alpha=0.3,
%             color="tab:green",
%         )
%         axes[0, col_idx].set_title(f"Feature {feat} | Observed vs Forecast")
%         axes[0, col_idx].legend()
%         axes[0, col_idx].grid(True)

%         axes[1, col_idx].plot(
%             time_axis[horizon_mask.numpy()],
%             trend_denorm[horizon_mask].numpy(),
%             color="tab:blue",
%         )
%         axes[1, col_idx].fill_between(
%             time_axis,
%             trend_lower_denorm.numpy(),
%             trend_upper_denorm.numpy(),
%             where=horizon_mask.numpy(),
%             alpha=0.3,
%             color="tab:blue",
%         )
%         axes[1, col_idx].set_title("Trend")
%         axes[1, col_idx].grid(True)

%         axes[2, col_idx].plot(
%             time_axis[horizon_mask.numpy()],
%             season_denorm[horizon_mask].numpy(),
%             color="tab:orange",
%         )
%         axes[2, col_idx].fill_between(
%             time_axis,
%             season_lower_denorm.numpy(),
%             season_upper_denorm.numpy(),
%             where=horizon_mask.numpy(),
%             alpha=0.3,
%             color="tab:orange",
%         )
%         axes[2, col_idx].set_title("Season")
%         axes[2, col_idx].grid(True)

%         axes[3, col_idx].plot(
%             time_axis[horizon_mask.numpy()],
%             jump_denorm[horizon_mask].numpy(),
%             color="tab:red",
%         )
%         axes[3, col_idx].fill_between(
%             time_axis,
%             jump_lower_denorm.numpy(),
%             jump_upper_denorm.numpy(),
%             where=horizon_mask.numpy(),
%             alpha=0.3,
%             color="tab:red",
%         )
%         axes[3, col_idx].set_title("Jump term")
%         axes[3, col_idx].grid(True)

%         axes[4, col_idx].plot(
%             time_axis[horizon_mask.numpy()],
%             error_denorm[horizon_mask].numpy(),
%             color="tab:purple",
%         )
%         axes[4, col_idx].fill_between(
%             time_axis,
%             error_lower_denorm.numpy(),
%             error_upper_denorm.numpy(),
%             where=horizon_mask.numpy(),
%             alpha=0.3,
%             color="tab:purple",
%         )
%         axes[4, col_idx].set_title("Error")
%         axes[4, col_idx].grid(True)

%         for row in range(nrows):
%             axes[row, col_idx].set_xlabel("Time step")
%             axes[row, col_idx].set_ylabel("Value")

%     plt.tight_layout()
%     plt.show()



% ================================================
% FILE: src/fits/modelling/FITSJ/model.py
% ================================================
% from dataclasses import dataclass

% import torch
% import torch.nn.functional as F
% from einops import reduce

% from fits.dataframes.dataset import ForecastingData
% from fits.modelling.FITSJ.transformer import Transformer, Decomposition
% from fits.modelling.framework import ForecastedData, ForecastingModel, ModelConfig


% @dataclass
% class FITSConfig(ModelConfig):
%     seq_len: int = 48
%     feature_size: int = 36
%     n_layer_enc: int = 4
%     n_layer_dec: int = 4
%     d_model: int = 64
%     n_heads: int = 4
%     mlp_hidden_times: int = 4
%     attn_pd: float = 0.0
%     resid_pd: float = 0.0
%     kernel_size: int | None = None
%     padding_size: int | None = None
%     lognormal_hucfg_t_sampling: bool = True  # otherwise, uniform
%     hucfg_num_steps: int = 500
%     first_differences: bool = True
%     conditional: bool = (
%         True  # use observed history values for conditioning during training/sampling
%     )

%     def fits_kwargs(self) -> dict[str, int | float | None]:
%         return {
%             "seq_length": self.seq_len,
%             "feature_size": self.feature_size,
%             "n_layer_enc": self.n_layer_enc,
%             "n_layer_dec": self.n_layer_dec,
%             "d_model": self.d_model,
%             "n_heads": self.n_heads,
%             "mlp_hidden_times": self.mlp_hidden_times,
%             "attn_pd": self.attn_pd,
%             "resid_pd": self.resid_pd,
%             "kernel_size": self.kernel_size,
%             "padding_size": self.padding_size,
%         }


% class FITSModel(ForecastingModel):
%     def __init__(self, config: FITSConfig = FITSConfig()):
%         super().__init__(config)
%         self.config: FITSConfig = config
%         self.model = Transformer(
%             n_feat=config.feature_size,
%             n_channel=config.seq_len,
%             n_layer_enc=config.n_layer_enc,
%             n_layer_dec=config.n_layer_dec,
%             n_heads=config.n_heads,
%             attn_pdrop=config.attn_pd,
%             resid_pdrop=config.resid_pd,
%             mlp_hidden_times=config.mlp_hidden_times,
%             max_len=config.seq_len,
%             n_embd=config.d_model,
%             conv_params=[config.kernel_size, config.padding_size],
%         ).to(self.device)

%         self.time_scalar = 1000  ## scale 0-1 to 0-1000 for time embedding

%     def forward(self, batch: ForecastingData):
%         diffusion_batch, partial_mask, _, forecast_mask = self._adapt_batch(batch)
%         return self._train_loss(diffusion_batch, forecast_mask, partial_mask)

%     @torch.no_grad()
%     def evaluate(self, batch: ForecastingData, n_samples: int) -> ForecastedData:
%         self.eval()

%         diffusion_batch, partial_mask, base_levels, _ = self._adapt_batch(batch)
%         batch_size = diffusion_batch.size(0)

%         samples = []
%         for _ in range(n_samples):
%             generated = self.fast_sample_infill(
%                 shape=(batch_size, self.config.seq_len, self.config.feature_size),
%                 target=diffusion_batch,
%                 partial_mask=partial_mask,
%             )
%             generated = generated.to(
%                 diffusion_batch.device,
%                 dtype=diffusion_batch.dtype,
%             )
%             samples.append(generated)

%         stacked = torch.stack(samples, dim=1)
%         if self.config.first_differences:
%             stacked = self._restore_levels(stacked, base_levels)

%         return ForecastedData(
%             forecasted_data=stacked,
%             forecast_mask=batch.forecast_mask.to(self.device, dtype=torch.float32),
%             observed_data=batch.observed_data.to(self.device, dtype=torch.float32),
%             observed_mask=batch.observed_mask.to(self.device, dtype=torch.float32),
%             time_points=batch.time_points[..., 0].to(self.device, dtype=torch.float32),
%         )

%     @torch.no_grad()
%     def decomposition_evaluate(
%         self, batch: ForecastingData, n_samples: int
%     ) -> tuple[ForecastedData, Decomposition]:
%         self.eval()

%         diffusion_batch, partial_mask, base_levels, _ = self._adapt_batch(batch)
%         batch_size = diffusion_batch.size(0)

%         samples = []
%         trend_samples = []
%         season_samples = []
%         jump_samples = []
%         error_samples = []

%         t_input = torch.ones(
%             batch_size, device=diffusion_batch.device, dtype=diffusion_batch.dtype
%         )
%         t_input = t_input * self.time_scalar

%         for _ in range(n_samples):
%             generated = self.fast_sample_infill(
%                 shape=(batch_size, self.config.seq_len, self.config.feature_size),
%                 target=diffusion_batch,
%                 partial_mask=partial_mask,
%             )
%             generated = generated.to(
%                 diffusion_batch.device,
%                 dtype=diffusion_batch.dtype,
%             )
%             samples.append(generated)

%             _, decomposed = self._output(generated, t_input)
%             trend_samples.append(decomposed.trend)
%             season_samples.append(decomposed.season)
%             jump_samples.append(decomposed.jump_term)
%             error_samples.append(decomposed.error)

%         stacked = torch.stack(samples, dim=1)
%         trend = torch.stack(trend_samples, dim=1)
%         season = torch.stack(season_samples, dim=1)
%         jump_term = torch.stack(jump_samples, dim=1)
%         error = torch.stack(error_samples, dim=1)

%         if self.config.first_differences:
%             stacked = self._restore_levels(stacked, base_levels)
%             trend = self._restore_levels(trend, base_levels)

%             zero_levels = torch.zeros_like(base_levels)
%             season = self._restore_levels(season, zero_levels)
%             jump_term = self._restore_levels(jump_term, zero_levels)
%             error = self._restore_levels(error, zero_levels)

%         return (
%             ForecastedData(
%                 forecasted_data=stacked,
%                 forecast_mask=batch.forecast_mask.to(self.device, dtype=torch.float32),
%                 observed_data=batch.observed_data.to(self.device, dtype=torch.float32),
%                 observed_mask=batch.observed_mask.to(self.device, dtype=torch.float32),
%                 time_points=batch.time_points[..., 0].to(
%                     self.device, dtype=torch.float32
%                 ),
%             ),
%             Decomposition(trend=trend, season=season, jump_term=jump_term, error=error),
%         )

%     def _output(self, x: torch.Tensor, t: torch.Tensor):
%         return self.model(x, t, padding_masks=None)

%     def _train_loss(
%         self,
%         x_start: torch.Tensor,
%         loss_mask: torch.Tensor,
%         partial_mask: torch.Tensor | None = None,
%     ) -> torch.Tensor:
%         z0 = torch.randn_like(x_start)
%         if self.config.conditional and partial_mask is not None:
%             z0 = z0.clone()
%             z0[partial_mask] = x_start[partial_mask]
%         z1 = x_start

%         t = torch.rand(z0.shape[0], 1, 1, device=z0.device)
%         if self.config.lognormal_hucfg_t_sampling:
%             t = torch.sigmoid(torch.randn(z0.shape[0], 1, 1, device=z0.device))

%         z_t = t * z1 + (1.0 - t) * z0
%         target = z1 - z0

%         model_out, _ = self._output(z_t, t.squeeze() * self.time_scalar)
%         train_loss = F.mse_loss(model_out, target, reduction="none")

%         loss_mask = loss_mask.to(device=train_loss.device, dtype=train_loss.dtype)
%         train_loss = train_loss * loss_mask

%         loss_sum = reduce(train_loss, "b ... -> b", "sum")
%         mask_sum = reduce(loss_mask, "b ... -> b", "sum").clamp_min(1.0)
%         return (loss_sum / mask_sum).mean()

%     def fast_sample_infill(
%         self,
%         shape: tuple[int, int, int],
%         target: torch.Tensor,
%         partial_mask: torch.Tensor | None = None,
%         hucfg_Kscale: float = 0.03,
%     ) -> torch.Tensor:
%         z0 = torch.randn(shape, device=self.device)
%         if self.config.conditional and partial_mask is not None:
%             z0 = z0.clone()
%             z0[partial_mask] = target[partial_mask]
%         z1 = zt = z0

%         for step in range(self.config.hucfg_num_steps):
%             t = step / self.config.hucfg_num_steps
%             t = t**hucfg_Kscale

%             z0 = torch.randn(shape, device=self.device)
%             if self.config.conditional and partial_mask is not None:
%                 z0 = z0.clone()
%                 z0[partial_mask] = target[partial_mask]

%             target_t = target * t + z0 * (1 - t)
%             zt = z1 * t + z0 * (1 - t)
%             if partial_mask is not None:
%                 zt = zt.clone()
%                 zt[partial_mask] = target_t[partial_mask]

%             v, _ = self._output(
%                 zt, torch.tensor([t * self.time_scalar], device=self.device)
%             )

%             z1 = zt.clone() + (1 - t) * v

%             if self.config.first_differences:
%                 z1 = torch.clamp(z1, min=-2, max=2)
%                 # in first differences [-1, 1] -> [-2, 2]=[(-1) - (1), 1 - (-1)]
%             else:
%                 z1 = torch.clamp(z1, min=-1, max=1)

%         return z1

%     def _adapt_batch(
%         self, batch: ForecastingData
%     ) -> tuple[torch.Tensor, torch.Tensor, torch.Tensor, torch.Tensor]:
%         observed_data = batch.observed_data.to(dtype=torch.float32, device=self.device)
%         context_mask = (batch.observed_mask * (1 - batch.forecast_mask)).to(self.device)
%         partial_mask = context_mask.bool()
%         # 0 - to be generated
%         # 1 - known at generation
%         forecast_mask = batch.forecast_mask.to(device=self.device, dtype=torch.float32)
%         base_levels = observed_data[:, :1]

%         if not self.config.first_differences:
%             return observed_data, partial_mask, base_levels, forecast_mask

%         diff_data = self._first_differences(observed_data)
%         diff_partial_mask = self._first_difference_mask(partial_mask)

%         return diff_data, diff_partial_mask, base_levels, forecast_mask

%     @staticmethod
%     def _first_differences(data: torch.Tensor) -> torch.Tensor:
%         diffs = torch.zeros_like(data)
%         diffs[:, 1:] = data[:, 1:] - data[:, :-1]
%         return diffs

%     @staticmethod
%     def _first_difference_mask(context_mask: torch.Tensor) -> torch.Tensor:
%         diff_mask = torch.zeros_like(context_mask)
%         diff_mask[:, 1:] = context_mask[:, 1:] & context_mask[:, :-1]
%         diff_mask[:, 0] = context_mask[:, 0]
%         return diff_mask

%     @staticmethod
%     def _restore_levels(
%         differences: torch.Tensor, base_levels: torch.Tensor
%     ) -> torch.Tensor:
%         base = base_levels.unsqueeze(1)
%         return base + differences.cumsum(dim=2)



% ================================================
% FILE: src/fits/modelling/FITSJ/transformer.py
% ================================================
% import os
% import math
% from dataclasses import dataclass

% import torch
% import numpy as np
% import torch.nn.functional as F
% from torch import nn
% from einops import rearrange, reduce, repeat

% from fits.modelling.FMTS.interpretable_diffusion.model_utils import (
%     AdaLayerNorm,
%     Conv_MLP,
%     GELU2,
%     RMSNorm,
%     Transpose,
% )


% ## hunote: our backbone network is most same as diffusion-TS. Diffusion-TS backbone has really good potential!!
% class TrendBlock(nn.Module):
%     """
%     Model trend of time series using the polynomial regressor.
%     """

%     def __init__(self, in_dim, out_dim, in_feat, out_feat, act):
%         super(TrendBlock, self).__init__()
%         trend_poly = 3
%         self.trend = nn.Sequential(
%             nn.Conv1d(
%                 in_channels=in_dim, out_channels=trend_poly, kernel_size=3, padding=1
%             ),
%             act,
%             Transpose(shape=(1, 2)),
%             nn.Conv1d(in_feat, out_feat, 3, stride=1, padding=1),
%         )

%         lin_space = torch.arange(1, out_dim + 1, 1) / (out_dim + 1)
%         self.poly_space = torch.stack(
%             [lin_space ** float(p + 1) for p in range(trend_poly)], dim=0
%         )

%     def forward(self, input):
%         b, c, h = input.shape
%         x = self.trend(input).transpose(1, 2)
%         trend_vals = torch.matmul(x.transpose(1, 2), self.poly_space.to(x.device))
%         trend_vals = trend_vals.transpose(1, 2)
%         return trend_vals


% class FourierLayer(nn.Module):
%     """
%     Model seasonality of time series using the inverse DFT.
%     """

%     def __init__(self, d_model, low_freq=1, factor=1):
%         super().__init__()
%         self.d_model = d_model
%         self.factor = factor
%         self.low_freq = low_freq

%     def forward(self, x):
%         """x: (b, t, d)"""
%         b, t, d = x.shape
%         x_freq = torch.fft.rfft(x, dim=1)

%         if t % 2 == 0:
%             x_freq = x_freq[:, self.low_freq : -1]
%             f = torch.fft.rfftfreq(t)[self.low_freq : -1]
%         else:
%             x_freq = x_freq[:, self.low_freq :]
%             f = torch.fft.rfftfreq(t)[self.low_freq :]

%         x_freq, index_tuple = self.topk_freq(x_freq)
%         f = repeat(f, "f -> b f d", b=x_freq.size(0), d=x_freq.size(2)).to(
%             x_freq.device
%         )
%         f = rearrange(f[index_tuple], "b f d -> b f () d").to(x_freq.device)
%         return self.extrapolate(x_freq, f, t)

%     def extrapolate(self, x_freq, f, t):
%         x_freq = torch.cat([x_freq, x_freq.conj()], dim=1)
%         f = torch.cat([f, -f], dim=1)
%         t = rearrange(torch.arange(t, dtype=torch.float), "t -> () () t ()").to(
%             x_freq.device
%         )

%         amp = rearrange(x_freq.abs(), "b f d -> b f () d")
%         phase = rearrange(x_freq.angle(), "b f d -> b f () d")
%         x_time = amp * torch.cos(2 * math.pi * f * t + phase)
%         return reduce(x_time, "b f t d -> b t d", "sum")

%     def topk_freq(self, x_freq):
%         length = x_freq.shape[1]
%         top_k = int(self.factor * math.log(length))
%         values, indices = torch.topk(
%             x_freq.abs(), top_k, dim=1, largest=True, sorted=True
%         )
%         mesh_a, mesh_b = torch.meshgrid(
%             torch.arange(x_freq.size(0)), torch.arange(x_freq.size(2)), indexing="ij"
%         )
%         index_tuple = (mesh_a.unsqueeze(1), indices, mesh_b.unsqueeze(1))
%         x_freq = x_freq[index_tuple]
%         return x_freq, index_tuple


% class SeasonBlock(nn.Module):
%     """
%     Model seasonality of time series using the Fourier series.
%     """

%     def __init__(self, in_dim, out_dim, factor=1):
%         super(SeasonBlock, self).__init__()
%         season_poly = factor * min(32, int(out_dim // 2))
%         self.season = nn.Conv1d(
%             in_channels=in_dim, out_channels=season_poly, kernel_size=1, padding=0
%         )
%         fourier_space = torch.arange(0, out_dim, 1) / out_dim
%         p1, p2 = (
%             (season_poly // 2, season_poly // 2)
%             if season_poly % 2 == 0
%             else (season_poly // 2, season_poly // 2 + 1)
%         )
%         s1 = torch.stack(
%             [torch.cos(2 * np.pi * p * fourier_space) for p in range(1, p1 + 1)], dim=0
%         )
%         s2 = torch.stack(
%             [torch.sin(2 * np.pi * p * fourier_space) for p in range(1, p2 + 1)], dim=0
%         )
%         self.poly_space = torch.cat([s1, s2])

%     def forward(self, input):
%         b, c, h = input.shape
%         x = self.season(input)
%         season_vals = torch.matmul(x.transpose(1, 2), self.poly_space.to(x.device))
%         season_vals = season_vals.transpose(1, 2)
%         return season_vals


% ### precompute_freqs_cis/reshape_for_broadcast/apply_rotary_emb are rope code adapted from LLama code
% def precompute_freqs_cis(dim: int, end: int, theta: float = 10000.0):
%     freqs = 1.0 / (theta ** (torch.arange(0, dim, 2)[: (dim // 2)].float() / dim))
%     t = torch.arange(end, device=freqs.device, dtype=torch.float32)
%     freqs = torch.outer(t, freqs)
%     freqs_cis = torch.polar(torch.ones_like(freqs), freqs)  # complex64
%     return freqs_cis


% def reshape_for_broadcast(freqs_cis: torch.Tensor, x: torch.Tensor):
%     ndim = x.ndim
%     assert 0 <= 1 < ndim
%     assert freqs_cis.shape == (x.shape[1], x.shape[-1])
%     shape = [d if i == 1 or i == ndim - 1 else 1 for i, d in enumerate(x.shape)]
%     return freqs_cis.view(*shape)


% def apply_rotary_emb(
%     xq: torch.Tensor,
%     xk: torch.Tensor,
%     freqs_cis: torch.Tensor,
% ):
%     xq_ = torch.view_as_complex(xq.float().reshape(*xq.shape[:-1], -1, 2))
%     xk_ = torch.view_as_complex(xk.float().reshape(*xk.shape[:-1], -1, 2))
%     freqs_cis = reshape_for_broadcast(freqs_cis, xq_)
%     xq_out = torch.view_as_real(xq_ * freqs_cis).flatten(3)
%     xk_out = torch.view_as_real(xk_ * freqs_cis).flatten(3)
%     return xq_out.type_as(xq), xk_out.type_as(xk)


% class FullAttention(nn.Module):
%     def __init__(
%         self,
%         n_embd,  # the embed dim
%         n_head,  # the number of heads
%         attn_pdrop=0.1,  # attention dropout prob
%         resid_pdrop=0.1,  # residual attention dropout prob
%         max_len=None,
%     ):
%         super().__init__()
%         assert n_embd % n_head == 0

%         self.key = nn.Linear(n_embd, n_embd)
%         self.query = nn.Linear(n_embd, n_embd)
%         self.value = nn.Linear(n_embd, n_embd)

%         self.attn_drop = nn.Dropout(attn_pdrop)
%         self.resid_drop = nn.Dropout(resid_pdrop)

%         self.proj = nn.Linear(n_embd, n_embd)
%         self.n_head = n_head

%         self.q_norm = RMSNorm(n_embd)
%         self.k_norm = RMSNorm(n_embd)

%         self.freqs_cis = precompute_freqs_cis(
%             n_embd // n_head,
%             max_len * 4,  ## if we use register, the overall length can be longer.
%             50000,  ##
%         )

%         self.regi_num = 128
%         self.register = nn.Parameter(torch.randn([1, self.regi_num, n_embd]))

%     def forward(self, x, mask=None):
%         # x = torch.cat([self.register.repeat(x.shape[0],1,1), x], 1)

%         B, T, C = x.size()
%         k = self.key(x)
%         q = self.query(x)
%         v = self.value(x)

%         k = self.k_norm(k) + 0.1 * k
%         q = self.q_norm(q) + 0.1 * q

%         k = k.view(B, T, self.n_head, C // self.n_head).transpose(
%             1, 2
%         )  # (B, nh, T, hs)
%         q = q.view(B, T, self.n_head, C // self.n_head).transpose(
%             1, 2
%         )  # (B, nh, T, hs)
%         v = v.view(B, T, self.n_head, C // self.n_head).transpose(
%             1, 2
%         )  # (B, nh, T, hs)

%         if int(os.environ.get("hucfg_attention_rope_use", "-1")) == 1:
%             freqs_cis = self.freqs_cis.cuda()[0:T]
%             q, k = apply_rotary_emb(
%                 q.permute(0, 2, 1, 3), k.permute(0, 2, 1, 3), freqs_cis=freqs_cis
%             )
%             q, k = q.permute(0, 2, 1, 3), k.permute(0, 2, 1, 3)

%         att = (q @ k.transpose(-2, -1)) * (1.0 / math.sqrt(k.size(-1)))  # (B, nh, T, T)

%         att = F.softmax(att, dim=-1)  # (B, nh, T, T)
%         # att = torch.sigmoid(att)
%         att = self.attn_drop(att)
%         y = att @ v  # (B, nh, T, T) x (B, nh, T, hs) -> (B, nh, T, hs)
%         y = (
%             y.transpose(1, 2).contiguous().view(B, T, C)
%         )  # re-assemble all head outputs side by side, (B, T, C)
%         att = att.mean(dim=1, keepdim=False)  # (B, T, T)

%         # output projection
%         y = self.resid_drop(self.proj(y))
%         # y = y[:,self.regi_num:,:]

%         return y, att


% class CrossAttention(nn.Module):
%     def __init__(
%         self,
%         n_embd,  # the embed dim
%         condition_embd,  # condition dim
%         n_head,  # the number of heads
%         attn_pdrop=0.1,  # attention dropout prob
%         resid_pdrop=0.1,  # residual attention dropout prob
%         max_len=None,
%     ):
%         super().__init__()
%         assert n_embd % n_head == 0
%         # key, query, value projections for all heads
%         self.key = nn.Linear(condition_embd, n_embd)
%         self.query = nn.Linear(n_embd, n_embd)
%         self.value = nn.Linear(condition_embd, n_embd)

%         # regularization
%         self.attn_drop = nn.Dropout(attn_pdrop)
%         self.resid_drop = nn.Dropout(resid_pdrop)
%         # output projection
%         self.proj = nn.Linear(n_embd, n_embd)
%         self.n_head = n_head

%         self.q_norm = RMSNorm(n_embd)
%         self.k_norm = RMSNorm(n_embd)

%         self.freqs_cis = precompute_freqs_cis(
%             n_embd // n_head,
%             max_len * 4,
%             50000,  ## hucfg913
%         )

%         self.regi_num = 128
%         self.register = nn.Parameter(torch.randn([1, self.regi_num, n_embd]))
%         self.register_2 = nn.Parameter(torch.randn([1, self.regi_num, n_embd]))

%     def forward(self, x, encoder_output, mask=None):

%         # x = torch.cat([self.register.repeat(x.shape[0],1,1), x], 1)

%         # encoder_output = torch.cat([self.register_2.repeat(x.shape[0],1,1), encoder_output], 1)

%         B, T, C = x.size()
%         B, T_E, _ = encoder_output.size()
%         # calculate query, key, values for all heads in batch and move head forward to be the batch dim
%         k = self.key(encoder_output)
%         q = self.query(x)

%         k = self.k_norm(k) + 0.1 * k  ## residual qk norm
%         q = self.q_norm(q) + 0.1 * q

%         k = k.view(B, T_E, self.n_head, C // self.n_head).transpose(1, 2)
%         q = q.view(B, T, self.n_head, C // self.n_head).transpose(
%             1, 2
%         )  # (B, nh, T, hs)

%         freqs_cis = self.freqs_cis.cuda()[0:T]
%         q, k = apply_rotary_emb(
%             q.permute(0, 2, 1, 3), k.permute(0, 2, 1, 3), freqs_cis=freqs_cis
%         )
%         q, k = q.permute(0, 2, 1, 3), k.permute(0, 2, 1, 3)

%         v = (
%             self.value(encoder_output)
%             .view(B, T_E, self.n_head, C // self.n_head)
%             .transpose(1, 2)
%         )  # (B, nh, T, hs)
%         att = (q @ k.transpose(-2, -1)) * (1.0 / math.sqrt(k.size(-1)))  # (B, nh, T, T)

%         att = F.softmax(att, dim=-1)  # (B, nh, T, T)
%         # att = torch.sigmoid(att)  ## sigmoid attention infact

%         att = self.attn_drop(att)
%         y = att @ v  # (B, nh, T, T) x (B, nh, T, hs) -> (B, nh, T, hs)
%         y = (
%             y.transpose(1, 2).contiguous().view(B, T, C)
%         )  # re-assemble all head outputs side by side, (B, T, C)
%         att = att.mean(dim=1, keepdim=False)  # (B, T, T)

%         y = self.resid_drop(self.proj(y))
%         # y = y[:,self.regi_num:,:]

%         return y, att


% class EncoderBlock(nn.Module):
%     """an unassuming Transformer block"""

%     def __init__(
%         self,
%         n_embd=1024,
%         n_head=16,
%         attn_pdrop=0.1,
%         resid_pdrop=0.1,
%         mlp_hidden_times=4,
%         activate="GELU",
%         max_len=None,
%     ):
%         super().__init__()

%         self.ln1 = AdaLayerNorm(n_embd)
%         self.ln2 = nn.LayerNorm(n_embd)
%         self.attn = FullAttention(
%             n_embd=n_embd,
%             n_head=n_head,
%             attn_pdrop=attn_pdrop,
%             resid_pdrop=resid_pdrop,
%             max_len=max_len,
%         )

%         assert activate in ["GELU", "GELU2"]
%         act = nn.GELU() if activate == "GELU" else GELU2()

%         self.mlp = nn.Sequential(
%             nn.Linear(n_embd, mlp_hidden_times * n_embd),
%             act,
%             nn.Linear(mlp_hidden_times * n_embd, n_embd),
%             nn.Dropout(resid_pdrop),
%         )

%     def forward(self, x, timestep, mask=None, label_emb=None):
%         a, att = self.attn(self.ln1(x, timestep, label_emb), mask=mask)
%         x = x + a
%         x = x + self.mlp(self.ln2(x))  # only one really use encoder_output
%         return x, att


% class Encoder(nn.Module):
%     def __init__(
%         self,
%         n_layer=14,
%         n_embd=1024,
%         n_head=16,
%         attn_pdrop=0.0,
%         resid_pdrop=0.0,
%         mlp_hidden_times=4,
%         block_activate="GELU",
%         max_len=None,
%     ):
%         super().__init__()

%         self.blocks = nn.Sequential(
%             *[
%                 EncoderBlock(
%                     n_embd=n_embd,
%                     n_head=n_head,
%                     attn_pdrop=attn_pdrop,
%                     resid_pdrop=resid_pdrop,
%                     mlp_hidden_times=mlp_hidden_times,
%                     activate=block_activate,
%                     max_len=max_len,
%                 )
%                 for _ in range(n_layer)
%             ]
%         )

%     def forward(self, input, t, padding_masks=None, label_emb=None):
%         x = input

%         for block_idx in range(len(self.blocks)):
%             x, _ = self.blocks[block_idx](x, t, mask=padding_masks, label_emb=label_emb)
%         return x


% class DecoderBlock(nn.Module):
%     """an unassuming Transformer block"""

%     def __init__(
%         self,
%         n_channel,
%         n_feat,
%         n_embd=1024,
%         n_head=16,
%         attn_pdrop=0.1,
%         resid_pdrop=0.1,
%         mlp_hidden_times=4,
%         activate="GELU",
%         condition_dim=1024,
%         max_len=None,
%     ):
%         super().__init__()

%         self.ln1 = AdaLayerNorm(n_embd)
%         self.ln2 = nn.LayerNorm(n_embd)
%         # self.ln2 = AdaLayerNorm(n_embd)

%         self.attn1 = FullAttention(
%             n_embd=n_embd,
%             n_head=n_head,
%             attn_pdrop=attn_pdrop,
%             resid_pdrop=resid_pdrop,
%             max_len=max_len,
%         )
%         self.attn2 = CrossAttention(
%             n_embd=n_embd,
%             condition_embd=condition_dim,
%             n_head=n_head,
%             attn_pdrop=attn_pdrop,
%             resid_pdrop=resid_pdrop,
%             max_len=max_len,
%         )

%         self.ln1_1 = AdaLayerNorm(n_embd)
%         # self.ln1_1 = nn.LayerNorm(n_embd)

%         assert activate in ["GELU", "GELU2"]
%         act = nn.GELU() if activate == "GELU" else GELU2()

%         self.trend = TrendBlock(n_channel, n_channel, n_embd, n_feat, act=act)
%         self.seasonal = FourierLayer(d_model=n_embd)

%         self.mlp = nn.Sequential(
%             nn.Linear(n_embd, mlp_hidden_times * n_embd),
%             act,
%             nn.Linear(mlp_hidden_times * n_embd, n_embd),
%             nn.Dropout(resid_pdrop),
%         )

%         self.proj = nn.Conv1d(n_channel, n_channel * 2, 1)
%         self.linear = nn.Linear(n_embd, n_feat)
%         self.jump_head = nn.Linear(n_embd, n_feat)

%     def forward(self, x, encoder_output, timestep, mask=None, label_emb=None):
%         a, att = self.attn1(self.ln1(x, timestep, label_emb), mask=mask)
%         x = x + a

%         a, att = self.attn2(self.ln1_1(x, timestep), encoder_output, mask=mask)
%         x = x + a
%         x1, x2 = self.proj(x).chunk(2, dim=1)
%         trend, season = self.trend(x1), self.seasonal(x2)
%         x = x + self.mlp(self.ln2(x))

%         m = torch.mean(x, dim=1, keepdim=True)
%         jump_score = self.jump_head(x)
%         return x - m, self.linear(m), trend, season, jump_score


% class Decoder(nn.Module):
%     def __init__(
%         self,
%         n_channel,
%         n_feat,
%         n_embd=1024,
%         n_head=16,
%         n_layer=10,
%         attn_pdrop=0.1,
%         resid_pdrop=0.1,
%         mlp_hidden_times=4,
%         block_activate="GELU",
%         condition_dim=512,
%         max_len=None,
%     ):
%         super().__init__()
%         self.d_model = n_embd
%         self.n_feat = n_feat
%         self.blocks = nn.Sequential(
%             *[
%                 DecoderBlock(
%                     n_feat=n_feat,
%                     n_channel=n_channel,
%                     n_embd=n_embd,
%                     n_head=n_head,
%                     attn_pdrop=attn_pdrop,
%                     resid_pdrop=resid_pdrop,
%                     mlp_hidden_times=mlp_hidden_times,
%                     activate=block_activate,
%                     condition_dim=condition_dim,
%                     max_len=max_len,
%                 )
%                 for _ in range(n_layer)
%             ]
%         )

%     def forward(self, x, t, enc, padding_masks=None, label_emb=None):
%         b, c, _ = x.shape
%         # att_weights = []
%         mean = []
%         season = torch.zeros((b, c, self.d_model), device=x.device)
%         trend = torch.zeros((b, c, self.n_feat), device=x.device)
%         jump_scores = torch.zeros((b, c, self.n_feat), device=x.device)
%         for block_idx in range(len(self.blocks)):
%             (
%                 x,
%                 residual_mean,
%                 residual_trend,
%                 residual_season,
%                 residual_jump,
%             ) = self.blocks[block_idx](
%                 x, enc, t, mask=padding_masks, label_emb=label_emb
%             )
%             season += residual_season
%             trend += residual_trend
%             jump_scores += residual_jump
%             mean.append(residual_mean)

%         mean = torch.cat(mean, dim=1)
%         return x, mean, trend, season, jump_scores


% @dataclass
% class Decomposition:
%     trend: torch.Tensor
%     season: torch.Tensor
%     jump_term: torch.Tensor
%     error: torch.Tensor


% class Transformer(nn.Module):
%     def __init__(
%         self,
%         n_feat,
%         n_channel,
%         n_layer_enc=5,
%         n_layer_dec=14,
%         n_embd=1024,
%         n_heads=16,
%         attn_pdrop=0.1,
%         resid_pdrop=0.1,
%         mlp_hidden_times=4,
%         block_activate="GELU",
%         max_len=2048,
%         conv_params=None,
%         **kwargs,
%     ):
%         super().__init__()
%         self.emb = Conv_MLP(n_feat, n_embd, resid_pdrop=resid_pdrop)
%         self.inverse = Conv_MLP(n_embd, n_feat, resid_pdrop=resid_pdrop)

%         if conv_params is None or conv_params[0] is None:
%             if n_feat < 32 and n_channel < 64:
%                 kernel_size, padding = 1, 0
%             else:
%                 kernel_size, padding = 5, 2
%         else:
%             kernel_size, padding = conv_params

%         self.combine_s = nn.Conv1d(
%             n_embd,
%             n_feat,
%             kernel_size=kernel_size,
%             stride=1,
%             padding=padding,
%             padding_mode="circular",
%             bias=False,
%         )
%         self.combine_m = nn.Conv1d(
%             n_layer_dec,
%             1,
%             kernel_size=1,
%             stride=1,
%             padding=0,
%             padding_mode="circular",
%             bias=False,
%         )
%         self.max_len = max_len
%         self.encoder = Encoder(
%             n_layer_enc,
%             n_embd,
%             n_heads,
%             attn_pdrop,
%             resid_pdrop,
%             mlp_hidden_times,
%             block_activate,
%             max_len=self.max_len,
%         )

%         self.decoder = Decoder(
%             n_channel,
%             n_feat,
%             n_embd,
%             n_heads,
%             n_layer_dec,
%             attn_pdrop,
%             resid_pdrop,
%             mlp_hidden_times,
%             block_activate,
%             condition_dim=n_embd,
%             max_len=self.max_len,
%         )
%         self.jump_kernel_size = 16
%         self.jump_filter = nn.Conv1d(
%             n_feat,
%             n_feat,
%             kernel_size=self.jump_kernel_size,
%             groups=n_feat,
%             bias=False,
%         )
%         self.jump_gain = nn.Parameter(torch.tensor(1.0))
%         self.jump_scale = nn.Parameter(torch.ones(1, 1, n_feat))
%         with torch.no_grad():
%             lags = torch.arange(self.jump_kernel_size, dtype=torch.float32)
%             decay = torch.exp(-lags / (self.jump_kernel_size / 3))
%             kernel = decay.view(1, 1, -1).repeat(n_feat, 1, 1)
%             self.jump_filter.weight.copy_(kernel)

%     def _causal_excitation(self, jump_scores: torch.Tensor) -> torch.Tensor:
%         scores = jump_scores.transpose(1, 2)
%         past_scores = F.pad(scores, (1, 0))[:, :, :-1]
%         padded = F.pad(past_scores, (self.jump_kernel_size - 1, 0))
%         excitation = self.jump_filter(padded)
%         return excitation.transpose(1, 2)

%     def forward(self, input, t, padding_masks=None, return_res=False):
%         emb = self.emb(input)

%         inp_enc = emb
%         enc_cond = self.encoder(inp_enc, t, padding_masks=padding_masks)

%         inp_dec = emb
%         output, mean, trend, season, jump_scores = self.decoder(
%             inp_dec, t, enc_cond, padding_masks=padding_masks
%         )

%         error = self.inverse(output)
%         error_m = torch.mean(error, dim=1, keepdim=True)
%         error = error - error_m

%         season = self.combine_s(season.transpose(1, 2)).transpose(1, 2)

%         excitation = self._causal_excitation(jump_scores)
%         excited_scores = jump_scores + self.jump_gain * excitation
%         jump = excited_scores * self.jump_scale
%         jump_m = torch.mean(jump, dim=1, keepdim=True)
%         jump = jump - jump_m

%         trend = self.combine_m(mean) + error_m + jump_m + trend

%         out = trend + season + jump + error
%         decomposed = Decomposition(trend, season, jump, error)

%         return out, decomposed

% \end{lstlisting}
